\documentclass[11pt]{amsart}
\usepackage[margin=1.1in]{geometry}
\usepackage{amsmath,amssymb,amsthm,booktabs,url,enumitem}
\usepackage[dvipsnames]{xcolor}\usepackage[pdftitle={Exploratory note: two-sided subset products for Williams numbers},pdfauthor={Raz Vakil},pdfkeywords={Williams numbers, Carmichael numbers, subset products}]{hyperref}
\hypersetup{colorlinks=true,linkcolor=blue!50!black,citecolor=blue!50!black,urlcolor=blue!50!black}
\providecommand{\lcm}{\operatorname{lcm}}
\title[Exploratory note: two-sided subset products]{Exploratory note to ``Exhaustive searches for Williams numbers'':\\ two-sided subset products}
\author{Raz Vakil}\address{Tel Aviv, Israel}\email{Raz.Vakil@gmail.com}\date{2 November 2026 (version 2; the onset section of version 1 is withdrawn)}
\begin{document}
\begin{abstract}
This unpublished note records one exploratory experiment that accompanied the computations of the main paper: a measurement of how the Erd\H{o}s--Alford--Granville--Pomerance subset-product construction degrades, in one specific implementation, when both $p-1$ and $p+1$ must be controlled. It is not a theorem and does not establish a general barrier; its limitations are listed in \S\ref{sec:limits}. Version 1 of this note also contained an Erd\H{o}s--Pomerance-style estimate of the size of the first Williams number; that section has been withdrawn (\S\ref{sec:withdrawn}) after a review showed that its upper endpoint was incompatible with the model's own pool-size calculation. Nothing in the main paper depends on this note. Code: \texttt{agrawal/} in the repository \url{https://github.com/razv2025/carmichael-lucas-carmichael}.
\end{abstract}
\maketitle
A \emph{Williams number} is a squarefree composite $n$ with $(p-1)\mid(n-1)$ and $(p+1)\mid(n+1)$ for every prime $p\mid n$; see the main paper for background, notation and references (cited here by the same keys).

\section{Two-sided subset products: a measurement}\label{sec:subset}
\emph{This note is heuristic.} 

In the subset-product method one fixes a smooth modulus $\Lambda$, harvests primes $p$ with $(p-1)\mid\Lambda$, and looks for a subset with product $\equiv1\pmod\Lambda$; the expected number of solutions is about $2^{\text{(pool size)}-\log_2\varphi(\Lambda)}$, and the method succeeds because the pool grows faster than $\log_2\varphi(\Lambda)$. For Williams numbers one needs simultaneously $(p+1)\mid\Lambda'$ and product $\equiv-1\pmod{\Lambda'}$, and the two moduli interact: an odd prime $\ell$ cannot divide both some $p_i-1$ and some $p_j+1$ (else $n\equiv1\equiv-1\pmod\ell$), so every small prime must be assigned to one side, and the $2$-adic branch is fixed. Define the \emph{margin} of a configuration as (number of usable primes) $-\log_2\varphi(\lcm\text{ of all required }p\mp1)$, so that the expected number of two-sided solutions is about $2^{\text{margin}}$.

We measured the margin over three families of configurations: random smooth moduli ($3{,}250$ configurations, all with margin $>-14$ searched exhaustively by meet-in-the-middle, no solution); pools of all primes $p\le X$ with both $p-1$ and $p+1$ $y$-smooth, with the side assignment optimised by hill-climbing (hundreds of restarts); and pools restricted to primes with few prime factors in $p\mp1$. The best margins were $-8.2$ (tiny three-prime systems), $-29.0/-30.6/-31.0/-36.9$ for $y=60/100/140/200$, and $-36.9$ for the sparse pools. The wall is flat or worsening in every direction we varied: each usable prime contributes one bit of subset freedom but about $2.5$ bits of modulus, because a prime that is useful on one side is almost never useful on the other. This is a measurement, not a theorem. The side-assignment of the small primes was formalised by Zhang \cite{Zhang15}, who splits the prime factors of the modulus into two groups serving $p-1$ and $p+1$ respectively, conjectures that the number of Williams numbers built from a modulus with $s$ prime factors grows only like $2^{2^{s^{1/2-o(1)}}}$ (against $2^{2^{s(1+o(1))}}$ for Carmichael numbers), and exhibits a set of $1{,}029$ primes expected to contain at least $2^{24}$ Williams numbers, none of which has been found. Our margins are a small-scale measurement consistent with Zhang's conjecture and with the ``intractable'' judgement of \cite{BFW21}.


\section{Limitations of the measurement}\label{sec:limits}
The figures above come from one implementation and should be read as such: side assignments are fixed per configuration and optimised by hill-climbing; the modulus is the lcm over the whole pool rather than over subsets; the greedy sub-pool handed to the meet-in-the-middle search is capped at 44 primes; meet-in-the-middle runs stop at a time limit without recording a completion flag; random-configuration harvesting stops after three million divisors, so pools may be truncated; six logged configurations above the stated margin threshold do not have an archived meet-in-the-middle run (direct checking confirms they are not solutions, but the log does not support the word ``exhaustively'' for them); and the best reported margin, $-8.2$, comes from a three-prime pool, where the margin counts all eight subsets although three-factor Williams numbers are impossible. The repository does not yet contain complete split logs, commands, hashes, parameters and completion records for the full experiment set. The measurement is consistent with Zhang's Conjecture~4 \cite{Zhang15} and with the judgement of \cite{BFW21}, but it does not prove anything about Erd\H{o}s--Alford--Granville--Pomerance-style constructions in general.

\section{Withdrawn: the onset estimate of version 1}\label{sec:withdrawn}
Version 1 of this note estimated, by an Erd\H{o}s--Pomerance-style count calibrated on twin-smooth prime counts, that the first Williams number lies near $10^{208}$--$10^{490}$ with about $80$--$110$ prime factors. A review of the model found that at the $10^{208}$ endpoint the optimiser selects $k=81$ prime factors and a smoothness bound $y\approx38.6$ with primes below $B\approx370$, where only $37$ primes meet the twin-smooth eligibility condition, while the model's fitted density multiplier (about $8.8$) predicts a pool of $548$; the endpoint is therefore impossible within the construction being modelled. Further defects: the four calibration samples imply per-prime factors varying from about $12.5$ to $3.75$ and do not justify a single constant; the residue factor uses $1/q$ where $q-1$ classes are relevant; the $2$-adic factor differs between the text and the code; the Agrawal variant double-counts a branch; the lower endpoint was obtained by forcing the expected count below $10^{18}$ not to exceed one, which is not a confidence bound; and the ``$561$ sanity check'' also produces an impossible pool density. The estimate is withdrawn and should not be cited. A corrected model would need bounded pool probabilities, correct residue and branch factors, explicit calibration uncertainty and recorded completion status; this has not been done.

\section*{Declaration of generative AI}
This note was drafted, and its code written, with the assistance of Claude (Anthropic); the review that found the defects above was produced with OpenAI Codex. The author is responsible for the content.

\begin{thebibliography}{9}
\bibitem{BFW21} R.~Baillie, A.~Fiori, S.~S.~Wagstaff Jr., Strengthening the Baillie--PSW primality test, \emph{Math. Comp.} 90 (2021), 1931--1955.
\bibitem{Zhang15} Z.~Zhang, Estimating the counts of Carmichael and Williams numbers with small multiple seeds, \emph{Math. Comp.} 84 (2015), 309--337.
\bibitem{Pinch07} R.~G.~E.~Pinch, Absolute quadratic pseudoprimes, Turku 2007.
\end{thebibliography}
\end{document}
